\section*{\S~7 Roleplay} 
\label{sec:roleplay}
\addcontentsline{toc}{section}{\S~7 Roleplay}

\begin{enumerate}[label=(\arabic*)]
    \item Die Roleplay-Regeln gelten ausschließlich während den Missionen. Für besondere Missionen können besondere Roleplay-Regeln verfasst werden. Diese müssen vorher mit der Leitung abgeklärt werden.   
    \item Grundsätzlich sollte immer ein seriöses Roleplay eingehalten werden, bei dem sich jeder entsprechend seiner Rolle verhält.    
    \item Es gelten dabei die bekannten Regeln:
    \begin{enumerate}[label=(\Alph*)]
        \item Fail RP: Es ist jederzeit Roleplay zu betreiben.
        \item Random Deathmatch (RDM): Das willkürliche Töten ist zu jeder Zeit verboten.
        \item Vehicle Deathmatch (VDM): Das Überfahren von Personen, von denen keine Gefahr ausgeht, ist ebenfalls verboten.
    \end{enumerate}    
    \item Vorgesetzte innerhalb der Mission sollten stets respektiert und deren Befehle ausgeführt werden. (\say{Gute Soldaten befolgen Befehle.} Zitat Chris, ehemaliger Leutnant)    
    \item Während der Mission hat derjenige mit der höchsten Funktion innerhalb der Mission unabhängig vom Dienstgrad das Kommando. Seinen Befehlen und Anweisungen ist Folge zu leisten.    
    \item Jeder (Unter)Offizier hat unmittelbar vor und nach der Mission zu jeder Zeit die Möglichkeit P(ermission)T(o)S(peak) auszurufen. PTS gilt, bis etwas anderes ausgerufen wird. Wer während PTS etwas sagen möchte, hat sich einmalig mit „PTS“ zu melden. Derjenige, der PTS ausruft, hat die Entscheidungsgewalt, ob und wann dem PTS stattgegeben wird. Das gleiche Recht steht dem Gefreiten vom Dienst (GvD) zu. PTS kann jederzeit von einem Vorgesetzten beendet werden. Innerhalb der Mission hat das dortige Führungspersonal dieses Recht inne. Mitglieder, die gegebenenfalls dennoch einen höheren Dienstgrad innehaben, haben Folge zu leisten.    
    \item Grundsätzlich steht der Spielspaß aller Mitglieder im Vordergrund. Diesen einem Mitglied zu verwehren oder zu vermindern gilt als schwerer Regelverstoß. Streitigkeiten oder Ähnliches können nach der Mission, gegebenenfalls im Beisein der Leitung, geklärt werden.
\end{enumerate}